Converting a hand-drawn circuit schematic into a form that an analog simulator can
run requires more than recognizing the drawing: it requires recovering the
\emph{electrical topology} --- which component terminals are connected to which.
Component and symbol recognition for engineering and circuit diagrams has advanced
considerably through detector-based pipelines, graph-enhanced postprocessing, and
domain datasets \cite{Thoma2021,Rachala2022,Shen2023,Schaefer2021,Elyan2020,Lee2020,Yang2024,Hemker2024,MorenoGarcia2024}.
Strong results are routinely reported for symbol detection, yet the conversion all
the way to a connected, simulatable netlist remains the harder and less solved
step, because wire recovery is sensitive to clutter, rotation, scan noise, and
faint traces, and because a single missed stroke leaves the recovered circuit
fragmented and therefore unsimulatable \cite{Rabby2019,Shen2023,Hemker2024}.

Wire recovery rests on classical document-analysis tools. Local adaptive
thresholding (Niblack, Sauvola, and contrast- or background-based refinements)
preserves thin foreground strokes on degraded documents far better than global
thresholding \cite{Otsu1979,Niblack1986,Sauvola2000,Gatos2006,Shafait2008,Su2013},
and skeletonization and line-segment detection are mature tools for tracing
elongated structure and preserving topology \cite{ZhangSuen1984,Lam1992,DudaHart1972,Gioi2010}.
These remain decisive when the objective is recovering sparse geometric structure
from weak, fragmented strokes rather than merely segmenting foreground.

Two gaps motivate the present method. First, recognition alone is insufficient:
work on piping, single-line, and circuit diagrams consistently finds that symbol
detection must be followed by domain-aware reasoning over lines, ports, and
connectivity to yield usable structure \cite{Elyan2020,Lee2020,Yang2024,Hemker2024,Rachala2022}.
General-purpose vision-language models are an attractive shortcut, but, as our
vision-language baseline in the Method validation shows, a strong VLM prompted to
read a schematic end to end recovers the components and gist while missing a large
fraction of the electrical topology ---
exactly the part that determines whether the circuit simulates. Explicit,
inspectable topology reasoning is therefore still needed. Second, the topology step
is hard to \emph{evaluate}: hand-drawn corpora annotate component boxes and wire
strokes, but not which terminals are electrically common, so there is no net-level
ground truth against which to measure or tune a join. This work addresses both: a
degree-budget topology join that explicitly recovers dropped connections, and a
synthetic evaluation harness that supplies net-level ground truth by construction
and was used to select the join.

The manuscript is written as a MethodsX article because the principal contribution
is methodological: a reproducible, component-guided pipeline whose core stage --
the Degree-Budget Topology Join -- is documented in enough detail to be inspected,
reused, and re-evaluated, together with the synthetic protocol used to validate it.
